%% kapa: exam=1 round=20 year=2009 period=1 form=A subject=civil_law subject_ko=민법 seq=34 number=34
%% source: Q-Net 감정평가사 기출문제 (cst00309) artlSeq=5024178 file=extracted/1008/5024178/2009년 제20회 감정평가사 1차시험 기출문제/1교시.hwp
%% conversion: kordoc 4.24.51 -> tools/gen1.py (auto); review=proofread-sonnet; flags=none
\begin{question}{34}
유치권에 관한 설명으로 옳지 \uline{않은} 것은? (다툼이 있으면 판례에 의함)
\begin{choices}
  \item 유치권자가 유치권을 행사하는 동안에는 피담보채권의 소멸시효가 진행하지 않는다.
  \item 수급인의 재료와 노력으로 신축된 건물에 대해서 특별한 사정이 없는 한 수급인은 유치권을 행사할 수 없다.
  \item 점유가 불법행위로 인한 경우에는 유치권이 성립하지 않는다.
  \item 채권의 변제기가 도래하지 않는 동안에는 유치권이 성립하지 않는다.
  \item 유치물이 경매된 경우 유치권자는 매수인(경락인)에게도 유치권을 주장할 수 있으나, 매수인에게 피담보채권의 변제를 청구할 수는 없다.
\end{choices}
\end{question}
\answer{1}
